%!TEX root = graduate-work.tex
% Добавьте ссылку на файлы с текстом работы
% Можно использовать команды:
%   \input или \include
% Пример:
%    \input{mainfiles/1-section} или \include{mainfiles/2-section}
% Команда \input позволяет включить текст файла без дополнительной обработки
% Команда \include при включении файла добавляет до него и после него команду
% перехода на новую страницу. Кроме того, она позволяет компилировать каждый файл
% в отдельности, что ускоряет сборку проекта.
% ВАЖНО: команда \include не поддерживает включение файлов, в которых уже содержится команда \include,
% т.е. не возможен рекурсивный вызов \include
\newcommand*{\Source}{
    %!TEX root = ../graduate-work.tex
\phantomsection
\section*{Введение} 
\addcontentsline{toc}{section}{Введение}

В этой части надо описать предметную область, задачу из которой вы будете решать, объяснить её актуальность (почему надо что-то делать сейчас?).
Здесь же стоит ввести определения понятий, которые вам понадобятся в постановке задачи.

Данная система основывается на $\mathbf{\mathbb {NP}}$-трудной проблеме ... Однако, у неё имеется серьёзный недостаток ...

В этой работе рассматривается ...  предложенное в 1994 коду~\cite{Ivanov}. В этой работе ...

В настоящей работе исследуются вопросы ...

В \S~1 даётся определение ...


Для некоторых классов приводятся нижние оценки ... (теорема~\ref{theorem1}).
    %!TEX root = ../graduate-work.tex

\section{Постановка задачи}

Здесь надо максимально формально описать суть задачи, которую потребуется решить, так, чтобы можно было потом понять, в какой степени полученное в результате работы решение ей соответствует. Текст главы должен быть написан в стиле технического задания, т.е. содержать как описание задачи, так и некоторый набор требований к решению

\textbf{Алгоритм Решения Задачи.}
\begin{itemize}
\item[1.] Описать Задачу
\item[2.] Декомпозировать задачу $a=mH\in\mathcal
C,\;wt(e')=\lfloor\frac{d-1}2\rfloor;$
\item[3.] Решить каждую подзадачу и объединить решение.
\end{itemize}
    %!TEX root = ../graduate-work.tex

\section{Обзор существующих решений}

Здесь надо рассмотреть все существующие решения поставленной задачи, но не просто пересказать, в чем там дело, а оценить степень их соответствия тем ограничениям, которые были сформулированы в постановке задачи.

Информацию о группах автоморфизмов других известных кодов можно найти в~\cite{FirstArticle}.
    %!TEX root = ../graduate-work.tex

\section{Исследование и построение решения задачи}
Здесь надо декомпозировать большую задачу из постановки на подзадачи и продолжать этот процесс, пока подзадачи не станут достаточно простыми, чтобы их можно было бы решить напрямую (например, поставив какой-то эксперимент или доказав теорему) или найти готовое решение.

\begin{theorem}\label{theorem1}

\begin{itemize}
    \item[1)] Гомоморфный образ группы
    \item[2)] Что во имя коммунизма
    \item[3)] Изоморфен фактор-группе
    \item[4)] По ядру гомоморфизма
\end{itemize}
\end{theorem}
\begin{proof}
    Доказательство оставляем читателю.
\end{proof}
    %!TEX root = ../graduate-work.tex

\section[Опи\-са\-ние практи\-чес\-кой час\-ти с переносами]{Описание практической части с переносами}

Если в рамках работы писался какой-то код, здесь должно быть его описание: выбранный язык и библиотеки и мотивы выбора, архитектура, схема функционирования, теоретическая сложность алгоритма, характеристики функционирования (скорость/память).
    %!TEX root = ../graduate-work.tex

\section{Заключение}

Здесь надо перечислить все результаты, полученные в ходе работы. Из текста должно быть понятно, в какой мере решена поставленная задача.

В каком-то смысле в заключении повторяются введе
}


% Информация о годе выполнения работы
\def\Year{%
    % 2006%
    2026%     % Текущий год
}

% Укажите тип работы
% Например:
%     Выпускная квалификационная работа,
%     Магистерская диссертация,
%     Курсовая работа, реферат и т.п.
\def\WorkType{%
    Численное решение модели реального делового цикла для закрытой экономики%
    % Магистерская диссертация%
    % Курсовая работа%
    % Реферат%
    %Дипломная работа%
}

% Название работы
%%%%%%%%%%% ВНИМАНИЕ! %%%%%%%%%%%%%%%%
% В МГУ ОНО ДОЛЖНО В ТОЧНОСТИ
% СООТВЕТСТВОВАТЬ ВЫПИСКЕ ИЗ ПРИКАЗА
% УТОЧНИТЕ НАЗВАНИЕ В УЧЕБНОЙ ЧАСТИ
\def\Title{%
    Численное решение модели реального делового цикла для закрытой экономики%
}


% Имя автора работы
\def\Author{%
    Козлов Кирилл Андреевич%
}

% Информация о научном руководителе
%% Фамилия Имя Отчество%
\def\SciAdvisor{%
    Богомолов Сергей Владимирович%
}
%% В формате: И.~О.~Фамилия%
\def\SciAdvisorShort{%
    С.~В.~Богомолов%
}
%% должность научного руководителя
\def\Position{%
    профессор%
    % доцент%
    % старший преподаватель%
    % преподаватель%
    % ассистент%
    % ведущий научный сотрудник%
    % старший научный сотрудник%
    % научный сотрудник%
    % младший научный сотрудник%
}
%% учёная степень научного руководителя
\def\AcademicDegree{%
    д.ф.-м.н.%
    % д.т.н.%
    % к.ф.-м.н.%
    % к.т.н.%
    % без степени%
}

% Информация об организации, в которой выполнена работа
%% Город
\def\Place{%
    Москва%
}
%% Университет
\def\Univer{%
    Московский государственный университет имени М.~В.~Ломоносова%
}
%% Факультет
\def\Faculty{%
    Факультет вычислительной математики и кибернетики%
}
%% Кафедра    
\def\Department{%
    Кафедра вычислительных методов%
}     

%%%% Переключите статус документа для отладки
%%%% В режиме draft документ собирается очень быстро
%%%% и выводится полезная информация о том
%%%% какие строки вылезают за границы документа, что удобно для борьбы с ними
\def\Status{%
    draft%
    % final%
}

%%%% Включает и выключает подпись <<С текстом работы ознакомлен>>
\def\EnableSign{%
    % true%
}
